\subsection{Objective and Input Image}

Task 3 requires extracting a sub-image that contains the complete character
row
\begin{equation}
  \texttt{HD44780A00}
  \label{eq:task3-target-string}
\end{equation}
The input is the $990\times367$ enhanced grayscale image selected in Task~1,
the 2--98\% percentile-stretched output saved under \texttt{task\_1/outputs/}.
The row of interest is the lower of the two text rows in the image; the upper
row is the logo \texttt{T12}.  The extraction must therefore not only crop the
correct vertical band but also decide the horizontal extent so that none of
the ten characters is cut off.

\subsection{Row Detection Method}

The row was located with projection profiles, which are computed entirely with
NumPy in the Task~3 script \path{task_3/task_3.py}.  The procedure has three
stages.

\subsubsection{Preprocessing for detection}

The enhanced image is first smoothed with a $5\times5$ box filter (the filter
implemented in Task~2) so that isolated bright speckles do not dominate the
profiles.  A binary mask is then obtained with Otsu's thresholding, whose
threshold $t^{*}$ maximises the between-class variance
\begin{equation}
  t^{*} =
  \arg\max_{t}\;
  \omega_0(t)\,\omega_1(t)\,
  \bigl(\mu_0(t)-\mu_1(t)\bigr)^{2}
  \label{eq:task3-otsu}
\end{equation}
where $\omega_0,\omega_1$ are the class probabilities and $\mu_0,\mu_1$ the
class means on either side of threshold $t$.  For the smoothed enhanced image
the Otsu threshold is $t^{*}=117$.

\subsubsection{Coarse row band}

The horizontal projection counts foreground pixels per image row,
\begin{equation}
  R(y)=\sum_{x} B(x,y) ,
  \label{eq:task3-row-projection}
\end{equation}
and the vertical projection counts foreground pixels per column,
\begin{equation}
  C(x)=\sum_{y} B(x,y) .
  \label{eq:task3-col-projection}
\end{equation}
Contiguous rows with $R(y)>25$ form candidate text bands.  The lowest band
below $y=185$ is the required row, giving the coarse box
$(x,y)\in[48,942)\times[198,336)$.

\subsubsection{Band-adaptive refinement}

The coarse box is looser than the true character extent because Otsu's mask
also marks part of the bright ring-like background texture.  The box is
therefore refined inside the coarse band: the band is thresholded again at its
80th intensity percentile (measured here as 183.6), which lies above the
textured background but below the character stroke cores.  Row spans shorter
than 40 pixels are discarded, and column spans separated by up to 30 pixels
are merged so that the spaces between characters do not split the row.  This
yields the refined box $(x,y)\in[62,941)\times[204,332)$.

\subsection{Extraction Result}

A safety margin of 10 pixels is added on every side and clamped to the image
bounds, giving the final extraction box
$(x,y)\in[52,951)\times[194,342)$, i.e.\ a sub-image of
$899\times148$ pixels.  \cref{fig:task3-region} shows the detected rectangle
overlaid on the full enhanced image, and \cref{fig:task3-subimage} shows the
extracted sub-image in both its original RGB form and the enhanced grayscale
form that continues into the later processing stages.

\begin{figure}[!htbp]
  \centering
  \includegraphics[width=\textwidth]{task3_extraction_region.png}
  \caption{Detected \texttt{HD44780A00} row with the extraction rectangle
  drawn in red.  The rectangle contains a 10-pixel safety margin around the
  refined character extent.}
  \label{fig:task3-region}
\end{figure}

\begin{figure}[!htbp]
  \centering
  \begin{subfigure}{\textwidth}
    \centering
    \includegraphics[width=\textwidth]{task3_subimage_rgb.png}
    \caption{Extracted sub-image, original RGB colour.}
    \label{fig:task3-subimage-rgb}
  \end{subfigure}

  \vspace{1ex}
  \begin{subfigure}{\textwidth}
    \centering
    \includegraphics[width=\textwidth]{task3_subimage.png}
    \caption{Extracted sub-image, enhanced grayscale (Task~1 output).}
    \label{fig:task3-subimage-gray}
  \end{subfigure}
  \caption{The $899\times148$ sub-image containing the complete
  \texttt{HD44780A00} row, extracted with the box
  $(x,y)\in[52,951)\times[194,342)$.}
  \label{fig:task3-subimage}
\end{figure}

\subsection{Verification}

The result was verified at the connected-component level.  Thresholding the
$5\times5$-smoothed crop at intensity 140 and labelling the components
recovers the ten characters as eight large stroke groups with areas between
$2{,}756$ and $8{,}933$ pixels; every other component is smaller than
$110$ pixels.  The labelled stroke groups span $x\in[60,941]$ and
$y\in[205,341]$, which lies entirely inside the final extraction box.  All
ten characters \texttt{H}, \texttt{D}, \texttt{4}, \texttt{4}, \texttt{7},
\texttt{8}, \texttt{0}, \texttt{A}, \texttt{0}, \texttt{0} are therefore
complete in the sub-image.  The \texttt{T12} logo, whose lowest stroke ends
at $y=181$, is not included because the crop starts at $y=194$.

\FloatBarrier
